% （随结果并入各问章后的）指定日结果与执行解释
\section{指定日期结果与执行解释}\label{sec:results}
\label{sec:reportdays}
题目要求的四个指定日期（3 月 20 日、6 月 21 日、9 月 23 日、12 月 21 日）结果均直接取自对应分支的完整全年连续运行，不为报告日重置库存。表~\ref{tab:days-q2}--\ref{tab:days-q4-3} 汇总四分支主方法在各指定日的全天购电量、实际总费用、紧急购电量与首末库存；每日 10 分钟粒度的完整购电与四小时充放电结果按题面表 1、表 2 格式存于对应 result 工作簿。

\begin{table}[!htbp]
\centering\small
\caption{问2指定日摘要（同一次连续运行）}\label{tab:days-q2}
\begin{tabular}{lrrrrr}
\toprule
日期 & 计划量/kWh & 费用/元 & 紧急量/kWh & 初库存/kWh & 末库存/kWh \\
\midrule
03-20 & 59036.55 & 43616.60 & 1260.45 & 8999.03 & 3833.60 \\
06-21 & 39461.83 & 22944.28 & 156.48 & 3095.64 & 9192.02 \\
09-23 & 66361.21 & 42715.51 & 0.00 & 6918.73 & 6319.21 \\
12-21 & 96407.68 & 62664.80 & 0.00 & 8754.14 & 9194.93 \\
\bottomrule
\end{tabular}
\end{table}

\begin{table}[!htbp]
\centering\small
\caption{问3指定日摘要（同一次连续运行）}\label{tab:days-q3}
\begin{tabular}{lrrrrr}
\toprule
日期 & 计划量/kWh & 费用/元 & 紧急量/kWh & 初库存/kWh & 末库存/kWh \\
\midrule
03-20 & 60324.10 & 40284.28 & 349.18 & 9039.09 & 3832.64 \\
06-21 & 39252.56 & 22727.23 & 127.58 & 3278.57 & 9177.10 \\
09-23 & 67119.11 & 42904.50 & 0.00 & 6938.79 & 6820.89 \\
12-21 & 96587.28 & 62778.17 & 0.00 & 8663.83 & 9139.03 \\
\bottomrule
\end{tabular}
\end{table}

\begin{table}[!htbp]
\centering\small
\caption{问4-2指定日摘要（同一次连续运行）}\label{tab:days-q4-2}
\begin{tabular}{lrrrrr}
\toprule
日期 & 计划量/kWh & 费用/元 & 紧急量/kWh & 初库存/kWh & 末库存/kWh \\
\midrule
03-20 & 60215.61 & 44713.33 & 1150.22 & 6014.42 & 1692.82 \\
06-21 & 41171.97 & 20780.33 & 156.21 & 2938.23 & 10495.24 \\
09-23 & 66856.50 & 44656.71 & 0.00 & 6798.08 & 6230.43 \\
12-21 & 95725.74 & 73248.40 & 0.00 & 10446.97 & 10380.74 \\
\bottomrule
\end{tabular}
\end{table}

\begin{table}[!htbp]
\centering\small
\caption{问4-3指定日摘要（同一次连续运行）}\label{tab:days-q4-3}
\begin{tabular}{lrrrrr}
\toprule
日期 & 计划量/kWh & 费用/元 & 紧急量/kWh & 初库存/kWh & 末库存/kWh \\
\midrule
03-20 & 62763.23 & 41944.66 & 335.55 & 6950.92 & 3907.20 \\
06-21 & 34042.43 & 17905.85 & 138.99 & 2811.62 & 4007.62 \\
09-23 & 66952.21 & 44755.87 & 0.00 & 6938.85 & 6257.72 \\
12-21 & 96561.91 & 73468.51 & 0.00 & 9562.16 & 10080.37 \\
\bottomrule
\end{tabular}
\end{table}


\begin{figure}[!htbp]\centering
\safeincludegraphics[width=0.92\textwidth]{fig_f12_exec_day}
\caption{3 月 20 日问题三主方法执行轨迹：上图为净负荷、最终交付购电量与紧急购电；下图为储电量与动态保留水平。午后光伏骤降形成的缺口大部分由储能覆盖，库存触及保留水平后转为少量紧急补购——保留水平挡住了“把库存放空去省 5 倍电价”的短视动作，为傍晚更贵的缺口留存了库存}
\label{fig:fig_f12_exec_day}
\end{figure}

图~\ref{fig:fig_f12_exec_day} 以 3 月 20 日为例给出问题三的完整执行解释。三个观察：其一，交付购电曲线在 12:00 后整体上移——12:00 预报下调了午后光伏，重优化以 1.5 倍价调增，替代了原本 5 倍价的紧急补购路径（当日紧急购电量 349\,kWh，问题 2 同日为 1300\,kWh）；其二，储电量在午后缺口段沿保留水平放电，而保留水平自身随剩余高价段减少而阶梯下移，日末放开至下限附近——未来不再有贵段时留存库存没有价值；其三，充电集中在凌晨低价段与午间光伏富余段，与问题 1 的确定性策略同构，说明因果执行继承了最优解的经济结构。单日案例用于解释已固定控制规则的行为，不用于反向选择参数。
